% =============================================================================
% Presentación: 30 Years of Adaptive Neural Networks
%   Perceptron, Madaline, and Backpropagation (Widrow et al.)
% Presentado por: Carissa Bush, Vidya Srinivas, Po-Chun Huang, Ming Hung Chen
% =============================================================================
\documentclass[aspectratio=169,utf8,spanish]{beamer}

% ── Paquetes ────────────────────────────────────────────────────────────────
\usepackage[utf8]{inputenc}
\usepackage[T1]{fontenc}
\usepackage[spanish,es-noshorthands]{babel}
\usepackage{amsmath,amssymb,amsthm}
\usepackage{graphicx}
\usepackage{booktabs}
\usepackage{hyperref}
\usepackage{xcolor}

% ── Tema ────────────────────────────────────────────────────────────────────
\usetheme{Madrid}
\usecolortheme{default}

% ── Metadatos ───────────────────────────────────────────────────────────────
\title[30 Years of Adaptive Neural Networks]{
  30 Years of Adaptive Neural Networks:\\
  Perceptron, Madaline, and Backpropagation
}
\subtitle{Widrow \textit{et al.}}
\author{Carissa Bush, Vidya Srinivas, Po-Chun Huang, Ming Hung Chen}
\institute{}
\date{}

% ── Atajos ──────────────────────────────────────────────────────────────────
\newcommand{\R}{\mathbb{R}}
\newcommand{\E}{\mathbb{E}}
\DeclareMathOperator{\sgn}{sgn}
\DeclareMathOperator{\sgm}{sgm}

\begin{document}

% ── Portada ─────────────────────────────────────────────────────────────────
\begin{frame}
  \titlepage
\end{frame}

% ═══════════════════════════════════════════════════════════════════════════
\section{Introducción}
% ═══════════════════════════════════════════════════════════════════════════

\begin{frame}{Introducción}
  \begin{itemize}
    \item Los conceptos de aprendizaje de máquina (\textit{machine learning}) existen hace décadas.
    \item Muchas ideas fueron descubiertas, y redescubiertas independientemente.
    \item Muchas ideas fueron inspiradas biológicamente.
    \item Todo comenzó con conceptos simples.
  \end{itemize}
\end{frame}

% ═══════════════════════════════════════════════════════════════════════════
\section{Conceptos Fundamentales}
% ═══════════════════════════════════════════════════════════════════════════

\begin{frame}{El Combinador Lineal Adaptativo}
  \begin{itemize}
    \item Produce una suma ponderada de las entradas.
    \item Cada entrada tiene un peso asociado.
    \item Los pesos pueden ser discretos o continuos.
    \item Los pesos pueden ser adaptados.
  \end{itemize}
  \begin{center}
    $y = \sum_{i} w_i x_i = \mathbf{W}^T \mathbf{X}$
  \end{center}
\end{frame}

\begin{frame}{Clasificación Binaria y Separabilidad Lineal}
  \textbf{Clasificación Binaria}
  \begin{itemize}
    \item La tarea objetivo es clasificar patrones de entrada en dos grupos: positivo y negativo.
  \end{itemize}

  \textbf{Separabilidad Lineal}
  \begin{itemize}
    \item Un conjunto de patrones de entrada es \textbf{linealmente separable} si existe una línea (o hiperplano) que puede separar puntos con etiquetas positivas de puntos con etiquetas negativas.
  \end{itemize}
\end{frame}

\begin{frame}{Un Clasificador Lineal Adaptativo: Adaline}
  \begin{center}
    \begin{tabular}{c|c|c}
      Vector de entrada &
      \begin{tabular}{c}
        Combinador Lineal \\
        Adaptativo
      \end{tabular} &
      \begin{tabular}{c}
        Elemento Umbral \\
        Signum
      \end{tabular} \\
      & & $\downarrow$ \\
      & & $\pm 1$
    \end{tabular}
  \end{center}
  \begin{itemize}
    \item Realiza tareas de clasificación binaria.
    \item Es capaz de realizar fronteras de separación lineales.
    \item Estas fronteras clasifican perfectamente conjuntos de datos linealmente separables.
  \end{itemize}
\end{frame}

% ═══════════════════════════════════════════════════════════════════════════
\section{Clasificadores No Lineales}
% ═══════════════════════════════════════════════════════════════════════════

\begin{frame}{Realizando Fronteras de Separación No Lineales}
  \begin{itemize}
    \item Cuando los datos no son linealmente separables en el espacio original, se necesitan técnicas no lineales.
    \item Dos enfoques principales:
    \begin{enumerate}
      \item Mapas de características (\textit{feature maps}): proyectar a un espacio de mayor dimensión.
      \item Combinaciones no lineales de fronteras lineales (Madaline, redes feedforward).
    \end{enumerate}
  \end{itemize}
\end{frame}

\begin{frame}{Clasificadores No Lineales: Mapas de Características}
  \begin{itemize}
    \item Si los patrones de entrada no son separables en el espacio original, se mapean a un espacio de mayor dimensión donde sí son linealmente separables.
  \end{itemize}
  \begin{center}
    Vector de entrada $\rightarrow$ Mapa de Características $\rightarrow$ Adaline $\rightarrow \pm 1$
  \end{center}
  \begin{itemize}
    \item Sea $\phi$ nuestro mapa de características.
    \item Ejemplo: $\phi: \R \to \R^2$, donde $\phi(x) = (x, x^2)$.
  \end{itemize}
\end{frame}

\begin{frame}{Ejemplo de Mapa de Características}
  \textbf{Espacio original (1D):} puntos $1, 2, 3$ no separables linealmente.
  \bigskip

  \textbf{Espacio mapeado (2D):}
  \begin{itemize}
    \item $1 \to (1, 1)$
    \item $2 \to (2, 4)$
    \item $3 \to (3, 9)$
  \end{itemize}
  \bigskip

  Frontera lineal en el espacio mapeado: $x_1^2 - 4x_1 + 3.5 = 0$.
\end{frame}

\begin{frame}{Clasificadores No Lineales: Madaline I}
  \begin{itemize}
    \item Si los patrones no son separables en el espacio original, se usa una combinación no lineal de fronteras de separación lineales.
  \end{itemize}
  \begin{center}
    \begin{tabular}{c}
      Vector de entrada \\
      $\downarrow$ \\
      \begin{tabular}{cc}
        Adaline & Adaline \\
      \end{tabular} \\
      $\downarrow$ \\
      Elemento Lógico Fijo (No Lineal) \\
      $\downarrow$ \\
      $\pm 1$
    \end{tabular}
  \end{center}
\end{frame}

\begin{frame}{Clasificadores No Lineales: Redes Feedforward}
  \begin{itemize}
    \item Combina múltiples Adalines en capas y alimenta las salidas de una capa a las entradas de la siguiente.
    \item Diferentes capas pueden tener diferentes funciones de activación.
  \end{itemize}
\end{frame}

\begin{frame}{Clasificadores No Lineales: Adaptación}
  \begin{itemize}
    \item Se puede calcular la señal de error en el vector de salida.
    \item La adaptación de la capa de salida es posible.
    \item ¿Qué hacemos con la adaptación de las capas ocultas?
    \item $\to$ Retro-propagación del error (\textit{backpropagation}).
  \end{itemize}
\end{frame}

% ═══════════════════════════════════════════════════════════════════════════
\section{Adaptación: El Principio de Mínima Perturbación}
% ═══════════════════════════════════════════════════════════════════════════

\begin{frame}{Adaptación — Principio de Mínima Perturbación}
  \textbf{Reglas de Corrección de Error (EC)}
  \begin{itemize}
    \item Alteran los pesos de una red para corregir el error en la respuesta de salida al patrón de entrada presente.
  \end{itemize}

  \textbf{Reglas de Descenso más Empinado (SD)}
  \begin{itemize}
    \item Alteran los pesos de una red mediante descenso de gradiente.
    \item Reducen el error promediado sobre todos los patrones de entrada.
  \end{itemize}
\end{frame}

% ═══════════════════════════════════════════════════════════════════════════
\section{Reglas de Corrección de Error}
% ═══════════════════════════════════════════════════════════════════════════

\begin{frame}{Reglas de Corrección de Error: Elemento Umbral Único}
  \begin{itemize}
    \item Cada nuevo patrón de entrada inicia un nuevo ciclo de adaptación.
    \item El objetivo es actualizar los pesos para reducir el error.
    \item Las reglas lineales hacen correcciones directamente proporcionales al error.
    \item Las reglas no lineales hacen correcciones que no son directamente proporcionales al error.
  \end{itemize}
  \begin{center}
    Patrón de entrada $\to$ Adaline $\to$ Regla EC $\to$ Error $\to$ Pesos actualizados
  \end{center}
\end{frame}

\begin{frame}{EC Elemento Umbral Único: $\alpha$-LMS}
  \begin{itemize}
    \item La actualización de pesos es independiente de la magnitud del patrón de entrada (auto-normalizante).
    \item La elección de $\alpha$ controla la estabilidad y velocidad de convergencia.
  \end{itemize}
  \begin{align*}
    \mathbf{W}_{k+1} &= \mathbf{W}_k + \alpha \frac{\varepsilon_k \mathbf{X}_k}{\|\mathbf{X}_k\|^2} \\[4pt]
    \varepsilon_k &= d_k - \mathbf{W}_k^T \mathbf{X}_k \\[4pt]
    \Delta\varepsilon_k &= -\alpha\varepsilon_k
  \end{align*}
  \begin{itemize}
    \item Rango típico: $0.1 < \alpha < 1$.
    \item Pesos inicializados usualmente en $0$.
  \end{itemize}
\end{frame}

\begin{frame}{EC Elemento Umbral Único: Regla del Perceptrón}
  \begin{itemize}
    \item Suma o resta patrones de entrada a los pesos para corregir el error.
    \item Garantiza convergencia para cualquier conjunto de patrones linealmente separable.
    \item Sin garantías de convergencia para patrones no linealmente separables.
  \end{itemize}
  \begin{align*}
    \mathbf{W}_{k+1} &= \mathbf{W}_k + \alpha \frac{\tilde{\varepsilon}_k}{2} \mathbf{X}_k
  \end{align*}
  donde $\tilde{\varepsilon}_k = d_k - y_k$ es el error del cuantizador (error binario).
\end{frame}

\begin{frame}{EC Elemento Umbral Único: Reglas de May}
  \begin{itemize}
    \item Todas las respuestas deseadas son $+1$ o $-1$.
    \item Define una ``zona muerta'' (\textit{dead zone}) o margen alrededor de $0$, denotada por $\pm\gamma$.
  \end{itemize}

  \textbf{Regla de Adaptación por Incremento}
  \begin{itemize}
    \item Patrón fuera del margen: se adapta por la regla del Perceptrón.
    \item Patrón dentro del margen: se adapta por una variante normalizada del Perceptrón de incremento fijo.
  \end{itemize}

  \textbf{Regla de Adaptación por Relajación Modificada}
  \begin{itemize}
    \item Patrón fuera del margen y correctamente clasificado: no se modifica.
    \item Patrón dentro del margen o mal clasificado: variante normalizada del Perceptrón.
    \item Si $\gamma \to \infty$, se convierte en $\alpha$-LMS.
  \end{itemize}
\end{frame}

% ═══════════════════════════════════════════════════════════════════════════
\section{Reglas de Corrección de Error: Redes Multi-Elemento}
% ═══════════════════════════════════════════════════════════════════════════

\begin{frame}{Madaline Rule I (MRI)}
  \begin{itemize}
    \item Permite la adaptación de un elemento de la primera capa porque el elemento lógico es fijo.
    \item La segunda capa consiste en un único elemento lógico de umbral fijo (OR, AND, voto mayoritario, etc.).
  \end{itemize}
  \begin{center}
    Vector de entrada $\to$
    \begin{tabular}{c}
      Adaline \\
      Adaline \\
      $\vdots$
    \end{tabular}
    $\to$ Lógica Fija $\to \pm 1$

    \small{\textit{Primera capa} \hspace{2cm} \textit{Segunda capa}}
  \end{center}
\end{frame}

\begin{frame}{Madaline Rule I: Adaptación}
  \begin{itemize}
    \item Los pesos de los Adalines se inicializan con valores aleatorios pequeños.
    \item El vector de pesos puede adaptarse mediante cualquiera de las reglas de corrección de error para un solo elemento:
    \begin{itemize}
      \item Invertir la salida del Adaline.
      \item $\alpha$-LMS.
      \item Perceptrón.
    \end{itemize}
    \item \textbf{Idea principal:} asignar responsabilidad al Adaline o Adalines que pueden asumirla más fácilmente.
  \end{itemize}
\end{frame}

\begin{frame}{Madaline Rule I: Características}
  \begin{itemize}
    \item Un asignador de trabajo (\textit{job assigner}) asigna la responsabilidad.
    \item El reparto de carga (\textit{load sharing}) es importante.
    \item La secuencia de presentación de patrones debe ser aleatoria.
    \item El proceso de adaptación puede atascarse en óptimos locales.
    \item Obedece el principio de mínima perturbación.
  \end{itemize}
\end{frame}

\begin{frame}{Madaline Rule II (MRII)}
  \begin{itemize}
    \item Extendido desde MRI.
    \item Usado para redes binarias multicapa.
    \item Los pesos se inicializan con valores aleatorios pequeños.
    \item Los patrones de entrenamiento se presentan en secuencia aleatoria.
    \item Puede atascarse en óptimos locales.
  \end{itemize}
\end{frame}

\begin{frame}{Madaline Rule II: Adaptación}
  \begin{itemize}
    \item El objetivo es reducir el error de Hamming.
    \item Si la red produce un error, adaptar la primera capa:
    \begin{enumerate}
      \item \textbf{Adaptación de prueba:} invertir la salida binaria.
      \item Agregar una perturbación $\Delta s$ sin adaptación.
      \item Si el error de salida se reduce, remover $\Delta s$ y adaptar los Adalines seleccionados mediante $\alpha$-LMS.
      \item Si el error no se reduce, no hay adaptación de pesos.
    \end{enumerate}
    \item Agotar todos los Adalines en la primera capa.
    \item Repetir para todas las capas.
  \end{itemize}
\end{frame}

% ═══════════════════════════════════════════════════════════════════════════
\section{Reglas de Descenso más Empinado}
% ═══════════════════════════════════════════════════════════════════════════

\begin{frame}{Reglas de Descenso más Empinado: Motivación}
  \begin{itemize}
    \item Las reglas de corrección de error corrigen una proporción del error para el patrón actual.
    \item Las reglas de descenso más empinado buscan minimizar el error cuadrático medio (MSE) sobre todos los patrones.
  \end{itemize}

  \begin{center}
    $\mathbf{W}_{k+1} = \mathbf{W}_k + \mu(-\nabla_k)$
  \end{center}

  donde $\mu$ controla la estabilidad y velocidad de convergencia, y $\nabla_k$ es el gradiente en $\mathbf{W} = \mathbf{W}_k$.
\end{frame}

\begin{frame}{Reglas de Descenso más Empinado: $\mu$-LMS}
  \begin{itemize}
    \item Se usa un gradiente crudo en lugar del gradiente real.
    \item El gradiente instantáneo se determina a partir de un solo patrón de entrada.
    \item $\mu$ es la constante de aprendizaje que determina estabilidad y tasa de convergencia.
  \end{itemize}
  \begin{align*}
    \hat{\nabla}_k &= \frac{\partial \varepsilon_k^2}{\partial \mathbf{W}_k} = -2\varepsilon_k \mathbf{X}_k \\
    \mathbf{W}_{k+1} &= \mathbf{W}_k + 2\mu\varepsilon_k \mathbf{X}_k
  \end{align*}
\end{frame}

\begin{frame}{Comparación: $\alpha$-LMS vs $\mu$-LMS}
  \begin{center}
    \begin{tabular}{ll}
      \toprule
      \textbf{$\alpha$-LMS} & \textbf{$\mu$-LMS} \\
      \midrule
      Gradiente instantáneo LMS & Gradiente instantáneo LMS \\
      Auto-normalizante & No auto-normalizante \\
      Más difícil de analizar & Más fácil de analizar \\
      Convergencia más rápida & Convergencia más lenta \\
      Puede converger a punto sesgado & Converge en media al mínimo (Wiener) \\
      \bottomrule
    \end{tabular}
  \end{center}
\end{frame}

\begin{frame}{Descenso más Empinado: Adaline con Sigmoide}
  \begin{itemize}
    \item Se extiende el Adaline para usar una función sigmoide en lugar de la función signum.
    \item $y_k = \sgm(s_k)$ donde $s_k = \mathbf{W}_k^T \mathbf{X}_k$.
  \end{itemize}
  \begin{align*}
    \tilde{\varepsilon}_k &= d_k - y_k \\
    \hat{\nabla}_k &= -2\tilde{\varepsilon}_k \sgm'(s_k) \mathbf{X}_k \\
    \mathbf{W}_{k+1} &= \mathbf{W}_k + 2\mu\tilde{\varepsilon}_k \sgm'(s_k) \mathbf{X}_k
  \end{align*}
\end{frame}

% ═══════════════════════════════════════════════════════════════════════════
\section{Retro-propagación (Backpropagation)}
% ═══════════════════════════════════════════════════════════════════════════

\begin{frame}{Retro-propagación (Backpropagation)}
  \textbf{Proceso:}
  \begin{enumerate}
    \item Entrada $\to$ Salida (propagación hacia adelante).
    \item Encontrar los errores en la salida.
    \item Propagar los efectos de los errores hacia atrás a través de la red para asociar una ``derivada del error cuadrático'' ($\delta$) con cada Adaline.
    \item Usar $\delta$ para determinar el gradiente.
    \item Actualizar los pesos basados en el gradiente.
    \item Repetir para todas las capas.
  \end{enumerate}
\end{frame}

\begin{frame}{Ejemplo de Backpropagation}
  \begin{itemize}
    \item Para una red feedforward de 3 capas:
  \end{itemize}
  \begin{align*}
    a^{(1)} &= \mathbf{X} \\
    a^{(2)} &= \sgm((W^{(1)})^T a^{(1)}) \\
    a^{(3)} &= \sgm((W^{(2)})^T a^{(2)}) \\
    a^{(4)} &= \sgm((W^{(3)})^T a^{(3)})
  \end{align*}
  \begin{align*}
    \delta_j^{(4)} &= a_j^{(4)} - y_j \quad \text{(error de salida)} \\
    \delta^{(3)} &= W^{(3)} \delta^{(4)} \;.\!*\; \sgm'(W^{(2)} a^{(2)}) \\
    \delta^{(2)} &= W^{(2)} \delta^{(3)} \;.\!*\; \sgm'(W^{(1)} a^{(1)})
  \end{align*}
\end{frame}

% ═══════════════════════════════════════════════════════════════════════════
\section{Madaline Rule III (MRIII)}
% ═══════════════════════════════════════════════════════════════════════════

\begin{frame}{Redes Multi-Elemento: Madaline Rule III}
  \begin{itemize}
    \item Motivación: el algoritmo de backpropagation requiere implementación precisa de la función sigmoide en hardware.
    \item MRIII usa una perturbación $\Delta s$ en lugar de la derivada de la sigmoide.
  \end{itemize}
  \begin{align*}
    \hat{\nabla}_k &= \frac{\partial \tilde{\varepsilon}_k^2}{\partial \mathbf{W}_k}
    \approx \frac{\Delta\tilde{\varepsilon}_k^2}{\Delta s} \mathbf{X}_k
    \approx 2\tilde{\varepsilon}_k \frac{\Delta\tilde{\varepsilon}_k}{\Delta s} \mathbf{X}_k
  \end{align*}
  \begin{itemize}
    \item MRIII es matemáticamente equivalente a backpropagation si $\Delta s$ es pequeño.
    \item MRIII es robusto incluso con implementación analógica.
  \end{itemize}
\end{frame}

\begin{frame}{MRIII vs Backpropagation}
  \begin{itemize}
    \item MRIII es esencialmente equivalente a Backpropagation:
    \begin{itemize}
      \item Mismos argumentos para cada uno de los elementos Adaline.
      \item El $\Delta s$ (perturbación) es pequeño.
      \item La adaptación se aplica a todos los elementos de la red a la vez.
    \end{itemize}
  \end{itemize}
\end{frame}

\begin{frame}{MRII vs MRIII}
  \begin{center}
    \begin{tabular}{p{0.45\textwidth}|p{0.45\textwidth}}
      \textbf{MRII} & \textbf{MRIII} \\
      \hline
      Discontinuo y no lineal & Todos los Adalines se adaptan \\
      No es posible usar gradientes instantáneos & Los que tienen sumas analógicas más cercanas a cero se adaptan más fuerte \\
      Problemas con mínimos locales & \\
    \end{tabular}
  \end{center}
\end{frame}

% ═══════════════════════════════════════════════════════════════════════════
\section{¿Qué se usa hoy?}
% ═══════════════════════════════════════════════════════════════════════════

\begin{frame}{¿Qué se sigue usando hoy?}
  \begin{itemize}
    \item Diferentes métodos son mejores en diferentes aplicaciones.
  \end{itemize}

  \textbf{Se usa menos:}
  \begin{itemize}
    \item Clasificación lineal.
    \item Elementos únicos.
    \item Una sola capa.
  \end{itemize}

  \textbf{Se usa más:}
  \begin{itemize}
    \item Multi-capa.
    \item Multi-elemento.
    \item Retro-propagación (\textit{backpropagation}).
  \end{itemize}
\end{frame}

% ═══════════════════════════════════════════════════════════════════════════
\section{Cierre}
% ═══════════════════════════════════════════════════════════════════════════

\begin{frame}{¿Preguntas?}
  \begin{center}
    \Large ¿Preguntas?
  \end{center}
\end{frame}

% ── Apéndice ────────────────────────────────────────────────────────────────
\appendix

\begin{frame}{Apéndice: Mapa de Características — Ejemplo Detallado}
  \textbf{Espacio original (1D):} puntos $1, 2, 3$.
  \bigskip

  \textbf{Mapa de características:} $\phi(x) = (x, x^2)$
  \begin{itemize}
    \item $1 \to (1, 1)$
    \item $2 \to (2, 4)$
    \item $3 \to (3, 9)$
  \end{itemize}

  \textbf{Frontera lineal en espacio mapeado:}
  $x_1^2 - 4x_1 + 3.5 = 0$
\end{frame}

\begin{frame}{Apéndice: $\mu$-LMS — Derivación}
  \begin{itemize}
    \item Sea $P^T \triangleq \E[d_k \mathbf{X}_k^T]$, $R \triangleq \E[\mathbf{X}_k \mathbf{X}_k^T]$.
    \item El gradiente del MSE: $\nabla_k = -2P + 2R\mathbf{W}_k$.
    \item Igualando a cero: $\mathbf{W}^* = R^{-1}P$ (vector de pesos de Wiener).
    \item $\mu$-LMS evita calcular $R^{-1}$ y $P$ usando el gradiente instantáneo.
    \item Para $0 < \mu < 1/\lambda_{\max}$, el algoritmo converge en media a $\mathbf{W}^*$.
  \end{itemize}
\end{frame}

\end{document}
